\documentclass[10pt,twocolumn]{article}
\usepackage[margin=0.75in]{geometry}
\usepackage{amsmath,amssymb}
\usepackage{booktabs}
\usepackage{graphicx}
\usepackage{caption}
\usepackage{hyperref}
\usepackage{setspace}
\usepackage{authblk}
\usepackage{microtype}
\usepackage{colortbl}
\usepackage{xcolor}

\setlength{\columnsep}{0.35in}
\setlength{\parskip}{2pt}
\setlength{\parindent}{12pt}
\captionsetup{font=small,labelfont=bf,skip=4pt}
\renewcommand{\baselinestretch}{1.02}

\title{\vspace{-1.2cm}\textbf{Does Double/Debiased Machine Learning Change the AJR Estimate?}\\
\textbf{A Coefficient-Comparison Test in a Small-Sample IV Design}}
\author[1]{Wenjing Gao}
\author[1]{Edith Jiang}
\affil[1]{EC610I -- Machine Learning in Economics, Wilfrid Laurier University}
\date{July 2026}

\begin{document}
\twocolumn[
\begin{@twocolumnfalse}
\maketitle
\begin{abstract}
\small
We ask a narrow, directly testable question: in Acemoglu, Johnson, and
Robinson's (2001, AJR) small-sample instrumental-variables design
($N=64$ former colonies), does replacing the linear first-stage and
outcome models with double/debiased machine learning (DML;
Chernozhukov et al., 2018) produce a coefficient that is
\emph{statistically different} from the conventional 2SLS estimate? We
estimate the institutional effect on log GDP per capita under three
nested control specifications, each with 2SLS and five DML nuisance
learners (Lasso, Ridge, Elastic Net, Random Forest, Gradient Boosting),
and report an approximate difference-in-coefficients test between
every DML estimate and its 2SLS counterpart from the identical
specification and sample. Across all 15 comparisons, no DML coefficient
is statistically distinguishable from 2SLS at the 5\% level. However,
the three regularized linear learners produce similar coefficients with
standard errors 25--45\% smaller than 2SLS. Tree-ensemble learners move
the point estimate substantially (in one specification, more than
doubling it) but their own standard errors inflate so much that even
this large a shift fails to reach significance -- evidence of estimator
instability, not a competing causal finding. These results suggest that
regularized linear DML can improve estimated precision without changing
the underlying causal conclusion, whereas tree-based nuisance learners
are unstable in this small sample.
\end{abstract}
\vspace{6pt}
\end{@twocolumnfalse}
]

\section{Introduction}
AJR (2001) instrument modern institutional quality with log settler
mortality to identify the causal effect of institutions on long-run GDP,
using 64 former European colonies. Their 2SLS estimate -- 0.93 in the
corrected 2012 data -- is a canonical result in development economics.
Rather than asking the broad question of whether ``machine learning
improves causal inference'' in the abstract, we ask a single, sharply
defined and directly testable question:

\begin{quote}
\small\itshape
If we replace the linear nuisance functions inside AJR's IV estimator
with machine learning, does the resulting coefficient differ, in a
formal statistical sense, from the original 2SLS estimate?
\end{quote}

This reframing matters because \emph{whether the coefficient differs}
and \emph{how precisely it is estimated} are two separate questions. A
large difference in point estimates is not, by itself, informative if
both estimates carry wide uncertainty; a small difference in point
estimates combined with a smaller standard error is informative even
when the two coefficients are not formally distinguishable. Our
contribution is to make this comparison explicit: for every
machine-learning estimator we consider, we report the point estimate, an
approximate equality test against 2SLS, and the percentage change in
standard error. This lets us describe, coefficient by coefficient,
whether -- and in what sense -- DML changes on the original 2SLS
approach.

Our focus on coefficient comparison distinguishes this paper from
studies that evaluate machine learning primarily by predictive
accuracy. In causal applications the object of interest is the
structural parameter $\beta_1$, not the fitted outcome, so predictive
performance alone is not sufficient evidence that a learner has
improved the analysis: a nuisance learner can fit $y$, $R$, or $M$
well and still produce a badly behaved $\hat\beta_1$ if it introduces
regularization bias into the residualization step. Accordingly, our
evaluation throughout emphasizes statistical equality of the
coefficient and changes in estimation precision, rather than
improvements in out-of-sample prediction.

\section{Data and Baseline}
Our sample is 64 former European colonies from the AJR (2012) reply
dataset. The endogenous regressor is protection against expropriation
risk (ICRG, 1985--1995, 0--10 scale). The instrument is log settler
mortality. The outcome is log GDP per capita.

We reproduce AJR (2012) Table 1B column 1 exactly: $\hat\beta_1 = 0.929$,
conventional SE $=0.156$, HC1-robust SE $=0.173$, first-stage $F=23.3$.
This no-controls estimate serves as a reference point throughout the
paper -- a benchmark estimate against which the fragility introduced by
adding controls can be judged -- though our formal equality tests
(Section 4) compare DML directly to the \emph{controlled} 2SLS estimate
from the same specification, not to this no-controls benchmark.

We estimate the causal effect under three nested control specifications
that span AJR's own robustness checks:
\begin{itemize}
\item \textbf{Spec 1 (latitude only, $N=64$):} the one control AJR
themselves treat as unambiguously pre-instrument.
\item \textbf{Spec 2 (AJR robustness controls, $N=62$):} latitude,
malaria prevalence (1994), and \% European descent (1975) -- AJR's own
robustness-check controls.
\item \textbf{Spec 3 (AJR + mediators, $N=56$):} Spec 2 plus log
population density, log trade, FDI, and life expectancy.
\end{itemize}
Specification 3 is best read as a sensitivity exercise rather than a
like-for-like extension of Specifications 1--2: because the added
variables may themselves be channels through which institutions affect
GDP, conditioning on them can shift the estimand toward a more
conditional institutional effect rather than AJR's total effect. Within
Specification 3, the comparison between DML and 2SLS is still valid
(both estimators target the same conditional parameter), but that
parameter should not be read as directly comparable to the total effect
recovered in Specifications 1--2. Each specification uses its own
complete-case sample, which is why $N$ varies across specifications.

\section{Method}
\subsection{DML in the IV Setting}
For each specification we estimate the causal effect two ways: we
estimate the original AJR specification using 2SLS and report
HC1-robust standard errors, and we apply DML in the partially linear IV
model of Chernozhukov et al.\ (2018).
Writing $y_i$ for log GDP, $R_i$ for institutional quality, $M_i$ for
log settler mortality, and $X_i$ for the control vector:
\begin{align}
y_i &= \beta_1 R_i + g_y(X_i) + \varepsilon_i, \\
R_i &= m_R(X_i) + \xi_i, \qquad M_i = m_M(X_i) + \eta_i.
\end{align}
Each nuisance function ($g_y$, $m_R$, $m_M$) is estimated by the same
machine-learning method on a training fold and residualized on a
held-out fold, using $K=5$-fold cross-fitting. The coefficient is
recovered from the cross-fitted residuals,
\begin{equation}
\hat\beta_1 = \frac{\sum_i \tilde M_i \tilde y_i}{\sum_i \tilde M_i \tilde R_i},
\end{equation}
Intuitively, DML first removes the variation explained by the control
variables from the outcome, the endogenous regressor, and the
instrument, using the fitted nuisance functions; the causal coefficient
is then estimated using only the residual variation that remains after
this residualization step. Cross-fitting ensures that the same
observations are never used both to train the nuisance models and to
form the residuals used in estimation, which limits the overfitting
bias that would otherwise result from a flexible learner fitting noise
in the training fold and passing that noise into $\hat\beta_1$.
with a heteroskedasticity-robust orthogonal-score standard error built
from the Neyman-orthogonal moment $\psi_i(\beta)=\tilde M_i(\tilde y_i -
\beta \tilde R_i)$. We compare this DML standard error against 2SLS's
HC1-robust standard error; both are heteroskedasticity-robust, but they
are constructed differently (one from a sandwich formula on
cross-fitted residuals, the other from the classical HC1 correction),
so the comparison is between two robust standard errors of different
provenance, not a like-for-like recomputation of the same object. We
use five nuisance learners: Lasso, Ridge, and Elastic Net (regularization
chosen by inner cross-validation), Random Forest (300 trees, depth 4),
and Gradient Boosting (150 trees, depth 3).

\subsection{The Coefficient-Comparison Test}
For every DML estimate we compute an approximate difference-in-coefficients
statistic against the 2SLS estimate from the identical specification and
sample:
\begin{equation}
z = \frac{\hat\beta_1^{\text{DML}} - \hat\beta_1^{\text{2SLS}}}
{\sqrt{\left(\text{SE}^{\text{DML}}\right)^2 + \left(\text{SE}^{\text{2SLS}}\right)^2}}.
\label{eq:ztest}
\end{equation}
This statistic is analogous in spirit to a Hausman-style comparison,
but it is not a formal Hausman test: it treats the two estimators'
sampling distributions as independent and does not estimate the
covariance between them, since 2SLS and DML are computed on the same
underlying sample and a tractable joint covariance is not available.
This is a simplifying approximation, not necessarily a conservative
one -- ignoring a nonzero covariance can move the test in either
direction depending on its sign. We report the resulting two-sided
$p$-value as our primary evidence on coefficient equality.
Confidence-interval overlap is reported alongside as a visual
diagnostic, not as an independent formal test; throughout the paper,
the formal comparison is the $p$-value from Equation~\ref{eq:ztest}.

\section{Results}

\subsection{Main Table: Coefficients and Coefficient-Comparison Tests}
Table~\ref{tab:main} is the central empirical result of the paper. For
every specification and every DML learner it reports the point
estimate, the standard error, the percentage change in SE relative to
2SLS, and the $p$-value from Equation~\ref{eq:ztest}.

\begin{table}[t]
\centering
\scriptsize
\begin{tabular}{lcccc}
\toprule
\textbf{Estimator} & $\hat\beta_1$ & \textbf{SE} & \textbf{SE $\Delta$} & \textbf{$p$} \\
\midrule
\multicolumn{5}{l}{\textit{Spec 1: Latitude only ($N{=}64$, $F{=}13.5$)}} \\
\midrule
2SLS               & 0.962 & 0.232 & --      & --    \\
DML Lasso          & 0.956 & 0.228 & $-$1.6\%  & 0.99 \\
DML Ridge          & 0.948 & 0.229 & $-$1.4\%  & 0.96 \\
DML Elastic Net    & 0.959 & 0.230 & $-$1.1\%  & 0.99 \\
DML Random Forest  & 0.910 & 0.247 & $+$6.5\%  & 0.88 \\
DML Gradient Boost.\ & 0.921 & 0.294 & $+$26.5\% & 0.91 \\
\midrule
\multicolumn{5}{l}{\textit{Spec 2: AJR robustness controls ($N{=}62$, $F{=}3.3$)}} \\
\midrule
2SLS               & 0.660 & 0.368 & --      & --    \\
DML Lasso          & 0.559 & 0.257 & $\mathbf{-30.2\%}$ & 0.82 \\
DML Ridge          & 0.609 & 0.275 & $\mathbf{-25.1\%}$ & 0.91 \\
DML Elastic Net    & 0.572 & 0.268 & $\mathbf{-27.2\%}$ & 0.85 \\
DML Random Forest  & 1.366 & 0.896 & $+$143.5\% & 0.47 \\
DML Gradient Boost.\ & 1.073 & 0.695 & $+$88.9\% & 0.60 \\
\midrule
\multicolumn{5}{l}{\textit{Spec 3: AJR + mediators ($N{=}56$, $F{=}2.8$)}} \\
\midrule
2SLS               & 0.450 & 0.336 & --      & --    \\
DML Lasso          & 0.422 & 0.236 & $\mathbf{-29.6\%}$ & 0.95 \\
DML Ridge          & 0.368 & 0.184 & $\mathbf{-45.1\%}$ & 0.83 \\
DML Elastic Net    & 0.413 & 0.227 & $\mathbf{-32.4\%}$ & 0.93 \\
DML Random Forest  & 0.694 & 0.341 & $+$1.6\%  & 0.61 \\
DML Gradient Boost.\ & 0.610 & 0.313 & $-$6.8\%  & 0.73 \\
\bottomrule
\end{tabular}
\caption{Coefficient comparison and approximate equality tests, DML
vs.\ 2SLS. SE is HC1-robust for 2SLS and orthogonal-score for DML. $p$
is the two-sided $p$-value from Equation~\ref{eq:ztest}. Every one of
the 15 DML-vs-2SLS pairs also has overlapping 95\% confidence intervals
(a secondary, visual check; see Figure~\ref{fig:main}). Bold SE changes
mark the $\geq$25\% reductions achieved by all three regularized linear
learners in Specs 2--3.}
\label{tab:main}
\end{table}

No coefficient in Table~\ref{tab:main} is statistically different from
its same-specification 2SLS counterpart: all 15 $p$-values exceed 0.46.
This means DML does not produce a formally different causal conclusion
than 2SLS at any specification we consider.

Where the two approaches diverge is precision, not the coefficient
itself. In Specifications 2 and 3 -- where the 2SLS first stage is
weakest (F = 3.3 and 2.8, both below the conventional threshold of 10)
-- the three regularized linear DML learners produce coefficients
statistically indistinguishable from 2SLS while cutting the standard
error by 25\% to 45\%. In Specification 3, for example, DML Ridge's
coefficient (0.368) is statistically indistinguishable from 2SLS's
(0.450), but its standard error (0.184) is roughly half of 2SLS's
(0.336). In Specification 1, where the instrument is already strong
(F = 13.5), the SE changes are much smaller ($-$1\% to $-$2\%).

Tree-ensemble learners move the point estimate substantially -- in
Specification 2, Random Forest's coefficient (1.366) is more than
double the 2SLS estimate (0.660) -- but the coefficient remains
statistically indistinguishable from 2SLS ($p=0.47$), because Random
Forest's own standard error (0.896) is more than double 2SLS's (0.368).
A large point-estimate shift combined with a large standard error is
not statistical evidence of a different effect.

\subsection{Main Figure}
Figure~\ref{fig:main} plots every estimator's point estimate and 95\%
confidence interval, grouped by specification, with 2SLS shown as a
black $\times$ in each panel. The pattern in Table~\ref{tab:main} is
directly visible: the three linear DML learners sit close to the 2SLS
estimate with visibly tighter intervals, while the two tree learners
are noticeably shifted with much wider intervals that still comfortably
contain the 2SLS point estimate.

\begin{figure}[t]
\centering
\includegraphics[width=\linewidth]{main_coefficient_comparison.png}
\caption{Estimated institutional effect: 2SLS and DML coefficients
across specifications. Points show coefficient estimates and horizontal
bars show 95\% confidence intervals. The black $\times$ is the
same-specification 2SLS estimate. The dashed line in each panel marks
that panel's own 2SLS coefficient, the benchmark used in the equality
test of Table~\ref{tab:main}.}
\label{fig:main}
\end{figure}

\subsection{Interpretation Across Specifications}
Read together, the three specifications reveal a consistent pattern
rather than three unrelated snapshots. When only latitude is controlled
for (Specification 1), the first-stage relationship remains relatively
strong (F = 13.5), and DML produces coefficients and standard errors
nearly identical to 2SLS -- there is little for a more flexible
nuisance learner to improve upon. As additional controls are introduced
in Specification 2, the classical first stage weakens considerably
(F = 3.3), and this is precisely where the regularized linear learners
begin to produce noticeably smaller standard errors while the
coefficient itself remains statistically indistinguishable from 2SLS.
Specification 3 introduces variables that may lie on the causal pathway
from institutions to income, so both 2SLS and DML there are best read
as estimating a more conditional effect (Section 2). Although the point
estimate declines further in this specification, the comparison between
DML and 2SLS \emph{within} it is unchanged from the pattern in
Specification 2: no learner produces a statistically different
coefficient, and the linear learners continue to deliver the same
coefficient more precisely. The precision advantage of linear DML,
in other words, tracks the weakness of the classical first stage
across specifications, not the specific set of controls used.

\section{Discussion}

\subsection{Linear Learners: Similar Coefficients, Smaller Standard Errors}
The standard-error reduction achieved by Lasso, Ridge, and Elastic Net
appears consistently across two independent specifications (2 and 3),
is largest exactly where the classical first stage is weakest, and
holds while the coefficient itself barely moves (differences of
0.03--0.09 in Specification 2, well within one standard error of either
estimator). The pattern is consistent with regularization stabilizing
nuisance-function estimation when controls are moderately collinear
with the instrument, as AJR's robustness controls are -- though we do
not separately test this mechanism against alternatives, and it should
be read as a plausible interpretation rather than an identified
mechanism. In Specification 1, where the classical first stage is
already strong, the linear DML learners offer almost no precision
advantage. This is economically intuitive rather than merely
descriptive: with a single control, the nuisance functions relating
latitude to the outcome, the treatment, and the instrument are close to
linear and easy to estimate with ordinary least squares. Replacing that
linear regression with a regularized machine-learning method provides
little additional flexibility to exploit, so DML has little room to
improve on 2SLS. The precision benefit becomes visible only once
additional, more collinear controls are introduced, where regularization
can reduce estimation noise in the nuisance functions while leaving the
recovered coefficient largely unchanged.

\subsection{Tree Learners: Large Shifts Without Statistical Support}
Random Forest's near-doubling of the coefficient in Specification 2
could be read, on its face, as evidence that a more flexible model
captures something 2SLS misses. The coefficient-comparison test cautions
against that reading: the difference is not statistically supported
($p=0.47$), because Random Forest's own standard error (0.896) is more
than double 2SLS's (0.368). At $N\approx 60$, a several-hundred-tree
ensemble has few effective observations per leaf, and we verified this
instability is not an artifact of our specific hyperparameters --
alternative depth and leaf-size specifications produce comparably large
standard errors. The large tree-based point-estimate shifts are better
interpreted as small-sample estimator instability than as evidence of
a different causal effect.

The tree learners illustrate a broader lesson that is easy to overlook
when evaluating machine learning by point estimates alone: a more
flexible learner is not necessarily a better learner in a causal
setting. Random Forest and Gradient Boosting produce substantially
different point estimates from 2SLS, but those differences are
accompanied by much larger standard errors -- in other words, the
additional flexibility increases variance faster than it reduces any
bias from nuisance-function misspecification. For a sample of roughly
sixty observations, the data simply do not contain enough information
to support a highly adaptive nuisance model with hundreds of trees;
the classical bias-variance tradeoff resolves here in favor of the
simpler, more constrained linear learners.

\subsection{In What Sense Does DML Improve on 2SLS?}
Two distinct answers, by learner class. For the regularized linear
learners, DML produces similar coefficients with materially smaller
estimated standard errors than 2SLS, in the specifications where 2SLS
is weakest. For the tree-ensemble learners, neither the coefficient nor
the precision improves: the point estimates move substantially, but the
accompanying standard errors are too large to treat those movements as
informative. Whether DML ``improves on'' 2SLS in this application
therefore depends entirely on which nuisance-learner class is used, and
answering that question requires the coefficient-by-coefficient
comparison reported here, rather than relying solely on visual
differences in point estimates.

\section{Limitations}
Three limitations deserve emphasis. First, our coefficient-comparison
test (Equation~\ref{eq:ztest}) does not estimate the covariance between
the 2SLS and DML estimators, even though both are computed on the same
underlying sample; a paired bootstrap that resamples both estimators
jointly would give a more exact test, and we do not implement that
here. Second, at $N=56$--64, both the DML sandwich standard error and
our normal-approximation test rely on asymptotic justifications that
warrant some caution at this sample size. The sample contains only
56--64 observations after complete-case deletion, and although DML
remains theoretically valid under the regularity conditions of
Chernozhukov et al.\ (2018), its finite-sample behavior depends heavily
on the quality of the underlying nuisance-function estimates -- which
is, in fact, the central empirical question this paper investigates.
Our findings should therefore be read as evidence from a particularly
challenging small-sample application, not as a general statement about
DML's performance at larger sample sizes. Third, the first stage itself
is weak in Specifications 2 and 3 (F below the conventional threshold
of 10), so point estimates from either estimator in those specifications
should be interpreted cautiously, independent of which method produced
them.

\section{Conclusion}
Across three specifications and five machine-learning nuisance
learners, no DML coefficient is statistically distinguishable from the
corresponding 2SLS estimate. DML does not change the substantive
conclusion of AJR's analysis. What differs by learner class is
precision: regularized linear DML learners produce similar coefficients
with materially smaller estimated standard errors than 2SLS, in the
specifications where the classical first stage is weakest. Tree-ensemble
learners add instability rather than value -- their point estimates
shift substantially, but the accompanying standard errors are too large
for those shifts to carry statistical weight.

More broadly, our results suggest that evaluating machine learning in
economics should not focus only on whether coefficients become larger
or smaller. Instead, researchers should ask whether the estimated
causal effect changes in a statistically meaningful way, and whether
the uncertainty surrounding that estimate improves. In the AJR
application, machine learning does not overturn the original economic
conclusion, but regularized linear DML suggests that a similar
causal estimate can sometimes be obtained more precisely -- which is a
narrower, but more defensible, claim than saying machine learning
``improves'' causal inference in general.

\begin{thebibliography}{9}\small
\bibitem{ajr2001} Acemoglu, D., Johnson, S., Robinson, J.A. (2001). The
Colonial Origins of Comparative Development. \emph{AER}, 91(5),
1369--1401.
\bibitem{ajr2012} Acemoglu, D., Johnson, S., Robinson, J.A. (2012). The
Colonial Origins of Comparative Development: Reply. \emph{AER}, 102(6),
3077--3110.
\bibitem{chern2018} Chernozhukov, V., et al.\ (2018). Double/Debiased
Machine Learning for Treatment and Structural Parameters.
\emph{Econometrics J.}, 21(1), C1--C68.
\bibitem{curtin1989} Curtin, P.D. (1989). \emph{Death by Migration}.
Cambridge University Press.
\bibitem{hausman1978} Hausman, J.A. (1978). Specification Tests in
Econometrics. \emph{Econometrica}, 46(6), 1251--1271.
\end{thebibliography}

\appendix
\section*{Appendix: Standard Errors by Specification and Estimator}
Figure~\ref{fig:se_comparison} shows the same standard errors reported
in Table~\ref{tab:main}, organized as a bar chart grouped by
specification and estimator, to make the SE pattern across the 18
estimates easier to scan visually alongside the coefficient-level
comparison in Figure~\ref{fig:main}.

\begin{figure}[h]
\centering
\includegraphics[width=\linewidth]{se_comparison.png}
\caption{HC1-robust (2SLS) and orthogonal-score (DML) standard errors
by specification and estimator. Bars correspond to the SE column of
Table~\ref{tab:main}.}
\label{fig:se_comparison}
\end{figure}

\end{document}